\documentclass[conference]{IEEEtran}
\usepackage{amsmath}
\usepackage{booktabs}
\usepackage{graphicx}
\usepackage{url}
\newcommand{\CARDIODIAGRunID}{20261107T202404Z}
\newcommand{\CARDIODIAGFigureDir}{figures}


\title{A Multi-Dataset Audit of CARDIODIAG: Nested Evaluation, Cohort Transfer, and Uncertainty on Public Cardiovascular Benchmarks}
\author{\IEEEauthorblockN{Akhil Preetham E, Ch Hari Krishna, and G Vennela}
\IEEEauthorblockA{Department of Computer Science and Engineering\\
Prasad V. Potluri Siddhartha Institute of Technology, Vijayawada, India\\
\textit{Author order, individual affiliations, contributions, and declarations require confirmation}}}

\begin{document}
\maketitle

\begin{abstract}
Public cardiovascular benchmarks can produce optimistic estimates when model selection, preprocessing, or cohort reuse is not controlled. We evaluated the existing CARDIODIAG pipeline on seven binary tasks with six models: Logistic Regression, K-Nearest Neighbors, Support Vector Machine, Random Forest, and the project’s unchanged OVEC and OSEC ensembles. Four UCI angiographic coronary artery disease (CAD) cohorts (920 records total) were evaluated separately, with three additional tasks for heart-failure death, myocardial-infarction lethal complications, and SPECT imaging diagnosis. We used nested stratified 3-fold cross-validation with training-only preprocessing and calibration partitions; the official SPECTF split was retained. All 72 directed transfers among the four compatible CAD feature schemas completed with source-only selection. The highest observed within-task ROC-AUC ranged from 0.679 to 0.911 across CAD cohorts and was 0.805 on the SPECTF held-out test set. OVEC and OSEC did not consistently outperform individual baselines. Cross-cohort ROC-AUC ranged from 0.238 to 0.903 across model-transfer pairs. Empirical coverage for nominal 90\% conformal sets ranged from 0.791 to 1.000, and source-versus-target feature-score AUROC varied from 0.357 to 0.934. These retrospective public-data experiments do not constitute independent clinical validation: the four UCI cohorts were represented in an earlier project composite, and participant-level overlap cannot be ruled out. Results support a reproducible benchmark audit, not clinical utility or deployment readiness.
\end{abstract}

\begin{IEEEkeywords}
cardiovascular prediction, coronary artery disease, nested cross-validation, dataset shift, calibration, conformal prediction, reproducibility
\end{IEEEkeywords}

\section{Introduction}
High internal discrimination is not sufficient evidence that a clinical prediction model is reliable at another site. Reuse of small public datasets, data-dependent preprocessing outside validation folds, and endpoint ambiguity can inflate estimates or weaken their clinical meaning. Reporting guidance and risk-of-bias frameworks therefore emphasize transparent data definitions, prediction timing, validation design, and model evaluation \cite{collins2024tripodai,moons2025probastai}. Nested resampling is needed when the same development data support hyperparameter selection and performance estimation \cite{varma2006bias,cawley2010overfitting}.

CARDIODIAG began as a feature-reduction and ensemble approach for heart-disease classification \cite{khan2026cardiodiag}. Previous project audits identified historical result artifacts, data-lineage concerns, and evaluation-boundary issues. This study preserves the established optimized voting classifier (OVEC) and stacking classifier (OSEC), then compares them with individual baselines using a leakage-controlled pipeline. The aims are to: (1) compare models across compatible and task-distinct public datasets without pooling outcomes; (2) assess discrimination, calibration, overfitting diagnostics, conformal coverage, selective risk, and feature-support alerts; and (3) measure source-only transfer among four schema-compatible CAD cohorts.

\section{Datasets and Study Design}
\subsection{Dataset selection and endpoints}
The UCI Heart Disease collection contains four historical cohorts, with 13 shared clinical predictors and the `num` angiographic disease field. We mapped `num=0` to negative and `num=1\ldots4` to positive, following the repository’s established binary rule \cite{ucihd}. The four component cohorts total 920 records. They are not independent replication datasets for this project: the same cohorts were represented in the previously examined Heart5 composite. Exact predictor-row matching found no cross-cohort exact duplicates, but no persistent participant identifiers were available to rule out transformed or near-duplicate records.

We also evaluated UCI Heart Failure Clinical Records (299 records), UCI Myocardial Infarction Complications (1,700 records), and UCI SPECTF Heart (267 records) \cite{ucihf519,ucimi579,ucispectf}. These were kept as separate tasks. For the heart-failure task, follow-up time was excluded as a predictor and death event was the target; follow-up horizons vary. For MI complications, the target was any lethal outcome and only admission-available variables were retained after excluding identifiers, later variables, and outcome proxies. SPECTF uses 44 SPECT-derived features and an overall normal/abnormal diagnosis; it is not a clinical CAD endpoint. The official 80-record training / 187-record test split was preserved, including 172 positive-coded and 15 negative-coded test records.

All source records remain separate. UCI dataset pages list CC BY 4.0; source URLs, archive checksums, class counts, missingness, and decisions for screened candidates are included in the accompanying dataset manifest and compatibility table. UCI Statlog Heart was excluded as a new cohort because the project’s exact matching audit found all 270 harmonized records match Cleveland \cite{ucistatlog}. Previously examined cohorts were not presented as new confirmation.

\begin{table*}[t]
\centering
\caption{Dataset characteristics. CAD cohorts share the UCI angiographic endpoint; heart failure, MI, and SPECTF are different tasks. Class counts are negative/positive.}
\label{tab:datasets}
\scriptsize
\begin{tabular}{lrrrrl}
\toprule
Dataset & $n$ & Predictors & Negative/positive & Missing cells & Task \\
\midrule
UCI Cleveland & 303 & 13 & 164/139 & 6 & Legacy angiographic CAD cohort; retrospective s... \\
UCI Hungary & 294 & 13 & 188/106 & 782 & Legacy angiographic CAD cohort; retrospective s... \\
UCI Switzerland & 123 & 13 & 8/115 & 396 & Legacy angiographic CAD cohort; retrospective s... \\
UCI VA Long Beach & 200 & 13 & 51/149 & 748 & Legacy angiographic CAD cohort; retrospective s... \\
UCI Heart Failure & 299 & 11 & 203/96 & 0 & Separate death-event task; varying follow-up ho... \\
UCI MI Complications & 1700 & 100 & 1429/271 & 15200 & Separate MI lethal-outcome task; admission vari... \\
UCI SPECTF Heart & 267 & 44 & 55/212 & 0 & Separate SPECT normal/abnormal task; official 8... \\
\bottomrule
\end{tabular}

\end{table*}

\subsection{Evaluation protocol}
For the six tabular cohorts evaluated with resampling, we used stratified three-fold outer validation and stratified three-fold inner model selection (seed 20261108). Within each outer training partition, 20\% was reserved for conformal and feature-support threshold calibration. Hyperparameter search used only the remaining development subset. Each outer test fold was held out until scoring. SPECTF used the official test partition; its training set was separated into development and calibration subsets only.

All imputation, scaling, categorical encoding, and outlier replacement were fitted within each pipeline’s training data. The existing `TrainingOnlyCleaner` replaces values outside training-fitted 1st/99th percentiles with training medians; numeric values are median-imputed and scaled; categorical values are most-frequent imputed and one-hot encoded. No SMOTE was applied. The corrected pipeline does not use a feature selector in this study. Consequently, no global feature selection or preprocessing was performed before splitting.

The individual model grids were logistic-regression $C\in\{0.1,1,10\}$; KNN $k\in\{3,5,7\}$ and uniform/distance weighting; SVM $C\in\{0.1,1,10\}$ and linear/RBF kernels; and random-forest square-root/all-feature subsampling and leaf size 1/3. OVEC and OSEC retained their existing definitions. OVEC uses the project’s soft-voting constituents (KNN, random forest, balanced logistic regression, and balanced SVC). OSEC retains its stacking architecture; its final estimator was tuned over $C\in\{0.1,1,10\}$.

\subsection{Metrics and uncertainty}
At a fixed probability threshold of 0.5, we report accuracy, precision, sensitivity, specificity, F1, ROC-AUC, and average precision (PR-AUC). Brier score and ten-bin equal-width expected calibration error (ECE), both macro-bin and sample-weighted, summarize calibration. Fold means and standard deviations describe nested evaluation. For pooled held-out predictions, 1,000 class-stratified percentile bootstrap replicates condition on observed class counts; these intervals do not account for patient clustering, cohort selection, or between-site variation. No hypothesis tests or superiority claims were prespecified.

Split-conformal sets use nonconformity $1-p(y)$ and nominal error level $\alpha=0.10$. Feature support is the mean distance to five nearest development observations after the fitted feature transformation; the alert cutoff is the 95th percentile of the source calibration scores. Cross-dataset domain AUROC distinguishes source calibration scores from target scores and is a feature-shift diagnostic, not an outcome-prediction metric. Confidence-ranked risk--coverage curves and the joint singleton-conformal / in-support accepted subset are descriptive. Conformal coverage guarantees require exchangeability and are not assumed to hold after transport \cite{angelopoulos2023conformal}.

\section{Results}
\subsection{Within-task model comparison}
All 42 dataset--model evaluations completed. Table~\ref{tab:auc} shows outer-fold ROC-AUC means and standard deviations; SPECTF is one official held-out test estimate and has no fold SD. The highest observed model by fold-mean ROC-AUC varied across tasks: random forest on Cleveland (0.911), OVEC on Hungary (0.890), logistic regression on Switzerland (0.772), random forest on VA Long Beach (0.679), OVEC on heart failure (0.752), random forest on MI complications (0.869), and OVEC on SPECTF (0.805). These maxima are descriptive comparisons, not unbiased estimates of a separately selected final model.

\begin{table*}[t]
\centering
\caption{ROC-AUC across all six candidates. Values are fold mean $\pm$ fold SD for nested evaluation. SPECTF reports the single official held-out test estimate.}
\label{tab:auc}
\scriptsize
\begin{tabular}{lrrrrrr}
\toprule
Dataset & LR & KNN & SVM & RF & OVEC & OSEC \\
\midrule
UCI Cleveland & 0.911 $\pm$ 0.031 & 0.889 $\pm$ 0.031 & 0.909 $\pm$ 0.039 & 0.911 $\pm$ 0.032 & 0.909 $\pm$ 0.033 & 0.903 $\pm$ 0.029 \\
UCI Hungary & 0.879 $\pm$ 0.040 & 0.876 $\pm$ 0.020 & 0.877 $\pm$ 0.043 & 0.880 $\pm$ 0.042 & 0.890 $\pm$ 0.033 & 0.876 $\pm$ 0.031 \\
UCI Switzerland & 0.772 $\pm$ 0.097 & 0.530 $\pm$ 0.120 & 0.569 $\pm$ 0.312 & 0.736 $\pm$ 0.096 & 0.732 $\pm$ 0.125 & 0.751 $\pm$ 0.183 \\
UCI VA Long Beach & 0.623 $\pm$ 0.026 & 0.641 $\pm$ 0.046 & 0.633 $\pm$ 0.036 & 0.679 $\pm$ 0.037 & 0.655 $\pm$ 0.041 & 0.612 $\pm$ 0.023 \\
UCI Heart Failure & 0.735 $\pm$ 0.039 & 0.661 $\pm$ 0.020 & 0.565 $\pm$ 0.324 & 0.739 $\pm$ 0.037 & 0.752 $\pm$ 0.027 & 0.736 $\pm$ 0.041 \\
UCI MI Complications & 0.854 $\pm$ 0.006 & 0.736 $\pm$ 0.023 & 0.846 $\pm$ 0.017 & 0.869 $\pm$ 0.008 & 0.866 $\pm$ 0.006 & 0.868 $\pm$ 0.009 \\
UCI SPECTF Heart & 0.758 & 0.729 & 0.788 & 0.796 & 0.805 & 0.797 \\
\bottomrule
\end{tabular}

\end{table*}

The complete thresholded and ranking results for each model appear in Tables~\ref{tab:cad_full} and~\ref{tab:other_full}; all metrics are also available as CSV. OVEC exceeded OSEC ROC-AUC on Cleveland, Hungary, heart failure, SPECTF, and VA Long Beach, while OSEC was higher on Switzerland and MI. Neither ensemble consistently exceeded all four individual baselines. Model selection results therefore do not support a general ensemble-superiority claim.

\begin{table*}[t]
\centering
\caption{Complete classification metrics for the four UCI CAD cohorts. Each value is outer-fold mean $\pm$ SD; PR-AUC should be read alongside cohort prevalence because precision depends on case mix.}
\label{tab:cad_full}
\scriptsize
\setlength{\tabcolsep}{2.2pt}
\resizebox{\textwidth}{!}{\begin{tabular}{llrrrrrrr}
\toprule
Dataset & Model & Acc. & Prec. & Sens. & Spec. & F1 & ROC-AUC & PR-AUC \\
\midrule
UCI Cleveland & K-Nearest Neighbors & 0.809 $\pm$ 0.029 & 0.819 $\pm$ 0.033 & 0.748 $\pm$ 0.033 & 0.860 $\pm$ 0.029 & 0.782 $\pm$ 0.032 & 0.889 $\pm$ 0.031 & 0.873 $\pm$ 0.027 \\
UCI Cleveland & Logistic Regression & 0.812 $\pm$ 0.026 & 0.797 $\pm$ 0.054 & 0.798 $\pm$ 0.034 & 0.823 $\pm$ 0.066 & 0.796 $\pm$ 0.018 & 0.911 $\pm$ 0.031 & 0.909 $\pm$ 0.030 \\
UCI Cleveland & OSEC & 0.825 $\pm$ 0.055 & 0.829 $\pm$ 0.091 & 0.791 $\pm$ 0.033 & 0.853 $\pm$ 0.099 & 0.808 $\pm$ 0.047 & 0.903 $\pm$ 0.029 & 0.901 $\pm$ 0.024 \\
UCI Cleveland & OVEC & 0.818 $\pm$ 0.035 & 0.808 $\pm$ 0.063 & 0.799 $\pm$ 0.024 & 0.835 $\pm$ 0.067 & 0.802 $\pm$ 0.029 & 0.909 $\pm$ 0.033 & 0.902 $\pm$ 0.031 \\
UCI Cleveland & Random Forest & 0.802 $\pm$ 0.030 & 0.795 $\pm$ 0.071 & 0.777 $\pm$ 0.035 & 0.823 $\pm$ 0.084 & 0.783 $\pm$ 0.016 & 0.911 $\pm$ 0.032 & 0.902 $\pm$ 0.031 \\
UCI Cleveland & Support Vector Machine & 0.828 $\pm$ 0.049 & 0.822 $\pm$ 0.073 & 0.806 $\pm$ 0.019 & 0.847 $\pm$ 0.078 & 0.813 $\pm$ 0.043 & 0.909 $\pm$ 0.039 & 0.908 $\pm$ 0.031 \\
UCI Hungary & K-Nearest Neighbors & 0.813 $\pm$ 0.021 & 0.790 $\pm$ 0.113 & 0.688 $\pm$ 0.106 & 0.882 $\pm$ 0.089 & 0.725 $\pm$ 0.018 & 0.876 $\pm$ 0.020 & 0.816 $\pm$ 0.039 \\
UCI Hungary & Logistic Regression & 0.793 $\pm$ 0.041 & 0.715 $\pm$ 0.086 & 0.735 $\pm$ 0.068 & 0.824 $\pm$ 0.091 & 0.720 $\pm$ 0.029 & 0.879 $\pm$ 0.040 & 0.822 $\pm$ 0.077 \\
UCI Hungary & OSEC & 0.806 $\pm$ 0.018 & 0.705 $\pm$ 0.034 & 0.802 $\pm$ 0.031 & 0.808 $\pm$ 0.044 & 0.749 $\pm$ 0.008 & 0.876 $\pm$ 0.031 & 0.821 $\pm$ 0.049 \\
UCI Hungary & OVEC & 0.799 $\pm$ 0.039 & 0.749 $\pm$ 0.087 & 0.688 $\pm$ 0.052 & 0.861 $\pm$ 0.077 & 0.713 $\pm$ 0.031 & 0.890 $\pm$ 0.033 & 0.837 $\pm$ 0.055 \\
UCI Hungary & Random Forest & 0.803 $\pm$ 0.046 & 0.748 $\pm$ 0.092 & 0.707 $\pm$ 0.093 & 0.856 $\pm$ 0.086 & 0.721 $\pm$ 0.053 & 0.880 $\pm$ 0.042 & 0.824 $\pm$ 0.062 \\
UCI Hungary & Support Vector Machine & 0.796 $\pm$ 0.047 & 0.759 $\pm$ 0.102 & 0.660 $\pm$ 0.052 & 0.872 $\pm$ 0.085 & 0.702 $\pm$ 0.043 & 0.877 $\pm$ 0.043 & 0.839 $\pm$ 0.057 \\
UCI Switzerland & K-Nearest Neighbors & 0.935 $\pm$ 0.014 & 0.935 $\pm$ 0.014 & 1 $\pm$ 0 & 0 $\pm$ 0 & 0.966 $\pm$ 0.007 & 0.530 $\pm$ 0.120 & 0.937 $\pm$ 0.026 \\
UCI Switzerland & Logistic Regression & 0.846 $\pm$ 0.061 & 0.938 $\pm$ 0.026 & 0.897 $\pm$ 0.092 & 0.167 $\pm$ 0.289 & 0.914 $\pm$ 0.038 & 0.772 $\pm$ 0.097 & 0.981 $\pm$ 0.011 \\
UCI Switzerland & OSEC & 0.902 $\pm$ 0.049 & 0.940 $\pm$ 0.030 & 0.956 $\pm$ 0.031 & 0.167 $\pm$ 0.289 & 0.948 $\pm$ 0.027 & 0.751 $\pm$ 0.183 & 0.976 $\pm$ 0.022 \\
UCI Switzerland & OVEC & 0.935 $\pm$ 0.014 & 0.935 $\pm$ 0.014 & 1 $\pm$ 0 & 0 $\pm$ 0 & 0.966 $\pm$ 0.007 & 0.732 $\pm$ 0.125 & 0.974 $\pm$ 0.017 \\
UCI Switzerland & Random Forest & 0.894 $\pm$ 0.037 & 0.941 $\pm$ 0.026 & 0.948 $\pm$ 0.068 & 0.167 $\pm$ 0.289 & 0.943 $\pm$ 0.022 & 0.736 $\pm$ 0.096 & 0.975 $\pm$ 0.015 \\
UCI Switzerland & Support Vector Machine & 0.935 $\pm$ 0.014 & 0.935 $\pm$ 0.014 & 1 $\pm$ 0 & 0 $\pm$ 0 & 0.966 $\pm$ 0.007 & 0.569 $\pm$ 0.312 & 0.945 $\pm$ 0.061 \\
UCI VA Long Beach & K-Nearest Neighbors & 0.715 $\pm$ 0.013 & 0.759 $\pm$ 0.008 & 0.906 $\pm$ 0.041 & 0.157 $\pm$ 0.068 & 0.826 $\pm$ 0.013 & 0.641 $\pm$ 0.046 & 0.815 $\pm$ 0.027 \\
UCI VA Long Beach & Logistic Regression & 0.610 $\pm$ 0.030 & 0.788 $\pm$ 0.008 & 0.651 $\pm$ 0.050 & 0.490 $\pm$ 0.034 & 0.713 $\pm$ 0.031 & 0.623 $\pm$ 0.026 & 0.812 $\pm$ 0.032 \\
UCI VA Long Beach & OSEC & 0.650 $\pm$ 0.061 & 0.845 $\pm$ 0.042 & 0.658 $\pm$ 0.143 & 0.627 $\pm$ 0.180 & 0.731 $\pm$ 0.080 & 0.612 $\pm$ 0.023 & 0.805 $\pm$ 0.027 \\
UCI VA Long Beach & OVEC & 0.720 $\pm$ 0.021 & 0.760 $\pm$ 0.012 & 0.913 $\pm$ 0.046 & 0.157 $\pm$ 0.090 & 0.829 $\pm$ 0.017 & 0.655 $\pm$ 0.041 & 0.818 $\pm$ 0.039 \\
UCI VA Long Beach & Random Forest & 0.730 $\pm$ 0.068 & 0.818 $\pm$ 0.017 & 0.819 $\pm$ 0.100 & 0.471 $\pm$ 0.059 & 0.817 $\pm$ 0.055 & 0.679 $\pm$ 0.037 & 0.838 $\pm$ 0.027 \\
UCI VA Long Beach & Support Vector Machine & 0.755 $\pm$ 0.008 & 0.753 $\pm$ 0.006 & 1 $\pm$ 0 & 0.039 $\pm$ 0.034 & 0.859 $\pm$ 0.004 & 0.633 $\pm$ 0.036 & 0.816 $\pm$ 0.026 \\
\bottomrule
\end{tabular}
}
\end{table*}

\begin{table*}[t]
\centering
\caption{Complete classification metrics for the task-distinct heart-failure, MI-complication, and SPECTF datasets. SPECTF has one official test split and no fold SD.}
\label{tab:other_full}
\scriptsize
\setlength{\tabcolsep}{2.2pt}
\resizebox{\textwidth}{!}{\begin{tabular}{llrrrrrrr}
\toprule
Dataset & Model & Acc. & Prec. & Sens. & Spec. & F1 & ROC-AUC & PR-AUC \\
\midrule
UCI Heart Failure & K-Nearest Neighbors & 0.685 $\pm$ 0.034 & 0.508 $\pm$ 0.125 & 0.229 $\pm$ 0.095 & 0.901 $\pm$ 0.009 & 0.314 $\pm$ 0.115 & 0.661 $\pm$ 0.020 & 0.457 $\pm$ 0.033 \\
UCI Heart Failure & Logistic Regression & 0.682 $\pm$ 0.046 & 0.508 $\pm$ 0.055 & 0.656 $\pm$ 0 & 0.695 $\pm$ 0.068 & 0.572 $\pm$ 0.036 & 0.735 $\pm$ 0.039 & 0.580 $\pm$ 0.073 \\
UCI Heart Failure & OSEC & 0.719 $\pm$ 0.033 & 0.554 $\pm$ 0.049 & 0.656 $\pm$ 0.054 & 0.749 $\pm$ 0.043 & 0.600 $\pm$ 0.044 & 0.736 $\pm$ 0.041 & 0.556 $\pm$ 0.018 \\
UCI Heart Failure & OVEC & 0.706 $\pm$ 0.015 & 0.559 $\pm$ 0.033 & 0.396 $\pm$ 0.036 & 0.852 $\pm$ 0.014 & 0.463 $\pm$ 0.033 & 0.752 $\pm$ 0.027 & 0.575 $\pm$ 0.024 \\
UCI Heart Failure & Random Forest & 0.696 $\pm$ 0.013 & 0.534 $\pm$ 0.019 & 0.396 $\pm$ 0.100 & 0.838 $\pm$ 0.038 & 0.450 $\pm$ 0.074 & 0.739 $\pm$ 0.037 & 0.544 $\pm$ 0.043 \\
UCI Heart Failure & Support Vector Machine & 0.679 $\pm$ 0.002 & 0 $\pm$ 0 & 0 $\pm$ 0 & 1 $\pm$ 0 & 0 $\pm$ 0 & 0.565 $\pm$ 0.324 & 0.455 $\pm$ 0.212 \\
UCI MI Complications & K-Nearest Neighbors & 0.851 $\pm$ 0.004 & 0.798 $\pm$ 0.189 & 0.096 $\pm$ 0.025 & 0.994 $\pm$ 0.006 & 0.169 $\pm$ 0.038 & 0.736 $\pm$ 0.023 & 0.415 $\pm$ 0.039 \\
UCI MI Complications & Logistic Regression & 0.807 $\pm$ 0.010 & 0.438 $\pm$ 0.017 & 0.727 $\pm$ 0.026 & 0.822 $\pm$ 0.017 & 0.546 $\pm$ 0.008 & 0.854 $\pm$ 0.006 & 0.565 $\pm$ 0.030 \\
UCI MI Complications & OSEC & 0.799 $\pm$ 0.007 & 0.428 $\pm$ 0.009 & 0.768 $\pm$ 0.021 & 0.805 $\pm$ 0.012 & 0.550 $\pm$ 0.003 & 0.868 $\pm$ 0.009 & 0.615 $\pm$ 0.014 \\
UCI MI Complications & OVEC & 0.874 $\pm$ 0.006 & 0.686 $\pm$ 0.016 & 0.380 $\pm$ 0.039 & 0.967 $\pm$ 0.001 & 0.489 $\pm$ 0.036 & 0.866 $\pm$ 0.006 & 0.597 $\pm$ 0.011 \\
UCI MI Complications & Random Forest & 0.869 $\pm$ 0.012 & 0.829 $\pm$ 0.064 & 0.236 $\pm$ 0.114 & 0.990 $\pm$ 0.009 & 0.356 $\pm$ 0.124 & 0.869 $\pm$ 0.008 & 0.625 $\pm$ 0.015 \\
UCI MI Complications & Support Vector Machine & 0.862 $\pm$ 0.006 & 0.638 $\pm$ 0.054 & 0.321 $\pm$ 0.040 & 0.965 $\pm$ 0.010 & 0.426 $\pm$ 0.036 & 0.846 $\pm$ 0.017 & 0.537 $\pm$ 0.023 \\
UCI SPECTF Heart & K-Nearest Neighbors & 0.529 & 0.988 & 0.494 & 0.933 & 0.659 & 0.729 & 0.964 \\
UCI SPECTF Heart & Logistic Regression & 0.738 & 0.949 & 0.756 & 0.533 & 0.841 & 0.758 & 0.975 \\
UCI SPECTF Heart & OSEC & 0.733 & 0.969 & 0.733 & 0.733 & 0.834 & 0.797 & 0.980 \\
UCI SPECTF Heart & OVEC & 0.749 & 0.970 & 0.750 & 0.733 & 0.846 & 0.805 & 0.980 \\
UCI SPECTF Heart & Random Forest & 0.743 & 0.956 & 0.756 & 0.600 & 0.844 & 0.796 & 0.979 \\
UCI SPECTF Heart & Support Vector Machine & 0.727 & 0.948 & 0.744 & 0.533 & 0.834 & 0.788 & 0.978 \\
\bottomrule
\end{tabular}
}
\end{table*}

\begin{figure*}[t]
\centering
\includegraphics[width=0.88\textwidth]{\CARDIODIAGFigureDir/model_performance_comparison.pdf}
\caption{Outer-fold mean ROC-AUC and PR-AUC with fold standard deviations. The SPECTF point is from its official test set and has no fold-based interval.}
\label{fig:performance}
\end{figure*}

\subsection{Overfitting, calibration, and selective diagnostics}
Resubstitution scores were often higher than outer-fold performance, particularly for KNN and random forest. For example, MI-complication training ROC-AUC was 1.000 for KNN and 1.000 (approximately) for random forest, while outer-fold means were 0.736 and 0.869. This is evidence of a training--evaluation gap in these fits, not proof that every high-score model is overfit. The full report includes model-specific training, selected inner-validation, outer-test values, and per-fold gaps.

Empirical coverage for nominal 90\% split-conformal sets varied from 0.791 to 1.000 across dataset--model groups; therefore the target was not uniformly met on these held-out samples. Average set sizes and singleton rates varied as well. Confidence-ranked selective coverage did not identify a universally reliable error-free region. The combined singleton-plus-support rule accepted between 0\% and 88.5\% of observations across task--model groups. One SPECTF model accepted none, while low-coverage zero-error subsets in other groups remain unstable and cannot certify safety.

The nearest-neighbor alert rate varied from 0.5\% to 37.4\% within held-out task predictions. Across cross-cohort transfers, source-versus-target feature-score AUROC ranged from 0.357 to 0.934, showing that feature-support separation depended strongly on source, target, and model. This distance score may detect some shifts in these samples, but does not provide a consistent rejection rule. Figure~\ref{fig:calibration} shows calibration behavior; the complete risk--coverage curves are included in the accompanying experimental report.

\begin{table*}[t]
\centering
\caption{Ranges across models for pooled Brier/ECE, mean empirical conformal coverage and set size, and combined singleton/support accepted coverage. These ranges summarize heterogeneous models; they are not model-selection recommendations.}
\label{tab:uncertainty}
\scriptsize
\begin{tabular}{lrrrrr}
\toprule
Dataset & Brier range & Weighted ECE range & Conformal coverage range & Set size range & Joint coverage range \\
\midrule
UCI Cleveland & 0.120--0.131 & 0.050--0.083 & 0.894--0.931 & 1.201--1.290 & 0.680--0.762 \\
UCI Heart Failure & 0.184--0.219 & 0.051--0.173 & 0.890--0.910 & 1.418--1.616 & 0.381--0.555 \\
UCI Hungary & 0.131--0.142 & 0.037--0.091 & 0.881--0.918 & 1.156--1.296 & 0.680--0.820 \\
UCI MI Complications & 0.095--0.147 & 0.017--0.212 & 0.905--0.922 & 1.086--1.347 & 0.624--0.885 \\
UCI SPECTF Heart & 0.168--0.329 & 0.233--0.455 & 0.791--1.000 & 1.128--2.000 & 0.000--0.866 \\
UCI Switzerland & 0.062--0.145 & 0.017--0.286 & 0.951--0.976 & 1.154--1.846 & 0.138--0.707 \\
UCI VA Long Beach & 0.182--0.236 & 0.047--0.216 & 0.900--0.970 & 1.595--1.843 & 0.155--0.380 \\
\bottomrule
\end{tabular}

\end{table*}

\begin{figure}[t]
\centering
\includegraphics[width=\columnwidth]{\CARDIODIAGFigureDir/calibration_curves.pdf}
\caption{Calibration curves from pooled held-out predictions. Bin counts and ECE variants are available in the result tables.}
\label{fig:calibration}
\end{figure}

\subsection{Source-only cross-cohort transfers}
All 72 directed model transfers across the four UCI CAD cohorts completed. Source-only tuning and calibration were used; target labels were used only for evaluation. Across model--direction pairs, target ROC-AUC ranged from 0.238 to 0.903. The average ROC-AUC across models varied by direction from 0.538 (Switzerland to VA Long Beach) to 0.886 (Cleveland to Hungary). The direction-level averages summarize model variability and should not obscure the full per-model matrix. These are retrospective stress tests because the component cohorts were part of the previously examined Heart5 lineage, and no patient identifier permits an independentness guarantee.

\begin{table*}[t]
\centering
\caption{Source-only transfer summary over the six models per direction. Mean and range summarize model scores on the same target cohort; they are descriptive. Domain OOD AUROC compares source calibration and target feature scores.}
\label{tab:transfer}
\scriptsize
\begin{tabular}{llrrrr}
\toprule
Source & Target & Mean AUC & Model range & Domain OOD AUROC & Target alert \\
\midrule
UCI Cleveland & UCI Hungary & 0.886 & 0.859--0.903 & 0.486 & 0.074 \\
UCI Cleveland & UCI Switzerland & 0.713 & 0.639--0.741 & 0.560 & 0.045 \\
UCI Cleveland & UCI VA Long Beach & 0.719 & 0.692--0.733 & 0.703 & 0.129 \\
UCI Hungary & UCI Cleveland & 0.843 & 0.820--0.873 & 0.896 & 0.421 \\
UCI Hungary & UCI Switzerland & 0.756 & 0.688--0.800 & 0.723 & 0.123 \\
UCI Hungary & UCI VA Long Beach & 0.713 & 0.689--0.739 & 0.727 & 0.160 \\
UCI Switzerland & UCI Cleveland & 0.635 & 0.534--0.773 & 0.874 & 0.612 \\
UCI Switzerland & UCI Hungary & 0.761 & 0.654--0.860 & 0.571 & 0.122 \\
UCI Switzerland & UCI VA Long Beach & 0.538 & 0.484--0.578 & 0.454 & 0.107 \\
UCI VA Long Beach & UCI Cleveland & 0.706 & 0.238--0.851 & 0.890 & 0.448 \\
UCI VA Long Beach & UCI Hungary & 0.813 & 0.702--0.868 & 0.772 & 0.235 \\
UCI VA Long Beach & UCI Switzerland & 0.662 & 0.610--0.711 & 0.586 & 0.046 \\
\bottomrule
\end{tabular}

\end{table*}

\begin{figure}[t]
\centering
\includegraphics[width=\columnwidth]{\CARDIODIAGFigureDir/cross_dataset_mean_auc_heatmap.pdf}
\caption{Mean target ROC-AUC over six source-only models for each CAD cohort direction. Source and target are historically examined UCI cohorts; diagonal cells are omitted.}
\label{fig:transfer}
\end{figure}

\subsection{Failed analyses}
Six learning-curve configurations failed, all for Switzerland. At the smallest requested training size, one or more stratified training subsets contained only one class; ROC-AUC learning curves were therefore not calculated for any model on that cohort. The remaining 36 task--model learning curves completed. All nested model evaluations and all cross-cohort transfers completed. Switzerland contains only eight negative records, so even the completed three-fold performance and OSEC stacking estimates have high uncertainty; the fixed five-fold internal stacking procedure generated minority-count warnings. The failures and warnings are preserved in the logs and failure CSV.

\section{Discussion}
This study broadens the CARDIODIAG audit while preserving the model definitions and keeping outcomes separate. In the compatible CAD cohorts, internal scores were strongest on Cleveland and Hungary and substantially weaker on VA Long Beach; Switzerland had only eight negative cases and yielded unstable fold estimates. The cross-cohort results varied sharply by direction and model. Even where the transfer AUC remained high, historical dataset reuse and unavailable participant linkage prevent treating these experiments as independent external validation.

OVEC and OSEC did not dominate the individual baselines consistently. Calibration and selective diagnostics were also mixed: nominal conformal coverage fell below 90\% in some model--task combinations, feature-support alerts varied widely, and low-coverage accepted subsets are not evidence of safe abstention. High apparent PR-AUC on imbalanced tasks must be interpreted against prevalence and the target definition; it does not by itself establish clinical value.

The additional public datasets illustrate why dataset count alone does not establish evidence strength. Heart-failure death and MI lethal complications are distinct outcomes with different prediction times. SPECTF is an imaging classification task and retains a severe class imbalance in its official test set. These tasks are useful for checking implementation behavior but cannot validate angiographic CAD prediction.

\subsection{Limitations}
The datasets are small, historical, and in some cases severely imbalanced or incomplete. Exact-row matching cannot detect near duplicates or transformations, while persistent participant identifiers are unavailable. The source tables do not establish a prospective index prediction time. Patient clustering, site sampling, and between-site variation are not represented by the class-stratified bootstrap intervals. Switzerland learning curves were not estimable under the requested training fractions, and OSEC's fixed internal stacking folds are poorly matched to such a small minority class. Conformal exchangeability is not established under dataset shift. No prospective, real-world, fairness, clinical-utility, or deployment study was performed.

\section{Conclusion}
The experiments provide a traceable comparison of individual classifiers and the existing CARDIODIAG ensembles across seven separately defined public tasks. No ensemble was consistently superior, cross-cohort performance varied substantially, and conformal/OOD/selective diagnostics did not produce a uniform reliability guarantee. These results support reproducible software and retrospective benchmark analysis. They do not establish independent clinical validation or readiness for use. A confirmatory study requires an authorized, untouched cohort with a prespecified CAD reference standard, prediction time, patient-level overlap controls, and a frozen operating protocol.

\section*{Declarations}
\textbf{Ethics and data governance:} Only public UCI repository datasets were used in this run. No institutional approval or exemption is claimed; the project team should obtain the appropriate institutional determination before submission.\\
\textbf{Funding and competing interests:} Confirmation pending.\\
\textbf{Author contributions:} Contributions and author order require confirmation from all listed authors.\\
\textbf{Data and code availability:} Raw records and row-level predictions are not included in the public release. Aggregate metrics, source checksums, preprocessing scripts, and figures are included in the project deliverables. Dataset attribution is required under the listed UCI licenses.\\
\textbf{Manuscript status:} Working draft requiring author, affiliation, declaration, and venue review.

\bibliographystyle{IEEEtran}
\bibliography{references}
\end{document}
